\section{Introduction}
\label{sec:introduction}

Preparing a prescribed quantum state is a basic component of quantum
simulation and variational quantum algorithms. General synthesis procedures
provide systematic constructions, but their circuit requirements grow rapidly
with the size of an arbitrary target~\cite{Mottonen2005,Plesch2011}.
Approximate preparation instead searches a restricted circuit family for a
state sufficiently close to the target. Its practical value depends on both
the resources required to execute the circuit and the classical work needed
to find it~\cite{Cerezo2021}.

Gate noise adds a second reason to control circuit size. Additional
parameterized gates can enlarge the set of accessible states, but executing
them also changes the noisy channel. The relation is not determined by gate
count alone. Different parameterizations of the same ideal state can have
different noise sensitivity, and redundant parameters can sometimes improve
the result of noisy reoptimization~\cite{Fontana2021}. Circuit reduction
therefore needs to be assessed through the fidelity of the resulting noisy
state, together with its construction cost.

Several approaches address parts of this problem. Noise-aware circuit
learning searches circuit structure and continuous parameters using a noisy
objective~\cite{Cincio2021}. Feedback-based refinement combines the measured
or simulated noisy output with derivatives of an ideal circuit
model~\cite{RobustState2024}. Parameter-efficient training alternates
optimization with pruning and regrowth~\cite{Sim2021PECT}. More recent
state-preparation work uses top-down sequential pruning, including selection
based on the cost of candidate gate removals~\cite{Ashfaq2026}.
Compilation-aware pruning also weighs the loss from omitting a rotation
against the noise associated with its routing overhead~\cite{Escofet2026}.
These methods establish circuit structure, parameter refinement, and noise
placement as coupled design choices.

Adjoint methods provide an efficient way to evaluate parameter sensitivities
in classical circuit simulation~\cite{JonesGacon2020}. Gate-local updates
inspired by Krotov optimal control use forward states and backward costates
to update parameters along the circuit~\cite{Goerz2019,Froehlich2026HybridKrotov}.
For the appropriate noiseless batch rule, this construction yields the
ordinary loss gradient. A sequential update uses changing forward states
while retaining backward information from the earlier iterate. These
ingredients can support circuit construction, but their use does not by
itself determine which gates to remove or how a sampled noisy refinement
should be interpreted.

Here, we study a specified combination of noiseless circuit growth,
target-dependent pruning, and noisy parameter refinement. At each visited
depth, sequential adjoint updates and limited-memory quasi-Newton refinement
construct an ideal circuit. We then evaluate the exact change in ideal
target fidelity caused by deleting each gate, remove gates below a fixed
threshold, and refine the surviving parameters using frozen noise histories.
Independent validation selects a depth before exact evaluation of the noisy
channel. The contribution is the construction and evaluation of this
workflow. We distinguish exact deletion scores and exact weighted history
compression from the implemented noisy update, for which we do not assume
an empirical-gradient identity or a monotonicity guarantee.

We evaluate the workflow on nine fixed six-qubit targets against three
publication-based simulator adaptations and two internal controls. Under
the synthetic logical-gate noise model, it gives the highest selected
fidelity on five targets and is Pareto-efficient in fidelity versus complete
construction and training CPU on seven. The real-valued targets provide the
strongest trade-offs. The complex Haar targets are more costly for the
workflow, and a separate native-gate comparison changes some rankings.
The results concern one optimizer seed and finite method-specific budgets;
they do not establish a general ranking of state-preparation algorithms.

We first define the state-preparation problem and gate-noise model. We then
describe the construction and selection procedure, followed by the numerical
protocol and the fidelity, runtime, and circuit-resource comparisons.
The discussion addresses the resulting scope and unresolved questions.
